\section{Requirements Elicitation and Analysis}\label{requirements}

Requirements were elicited from the intended project goal described in \cref{concept}: build a decentralized laboratory platform for distributed sensor ingestion, replicated state propagation, fault observation, and monitoring. The elicitation process therefore focused on the observable services the platform must provide to operators and observers, on the quality attributes expected from a distributed-system prototype, and on the practical constraints imposed by course evaluation and reproducible experimentation.

\subsection{Glossary}
\label{requirements-glossary}

\begin{itemize}
  \item \textbf{Node}: autonomous peer process participating in the cluster and executing the full runtime stack.
  \item \textbf{Peer}: another node known by a node through bootstrap, heartbeat exchange, or gossip dissemination.
  \item \textbf{Sensor reading}: atomic observation produced by a local sensor source, including sensor identifier, value, timestamp, and origin.
  \item \textbf{Replicated sensor state}: node-local view of the cluster-wide latest known sensor values after merge.
  \item \textbf{Membership view}: node-local knowledge about known peers and their liveness status.
  \item \textbf{Observer}: human user or external tool consuming read-only monitoring data.
  \item \textbf{Operator}: human user responsible for starting nodes, configuring experiments, and injecting or observing failures.
  \item \textbf{Convergence}: condition in which reachable replicas expose the same winning sensor values after enough communication.
  \item \textbf{Partition}: temporary communication disruption that prevents some nodes from exchanging messages with the rest of the cluster.
\end{itemize}

\subsection{Functional Requirements}
\label{functional-requirements}

\textbf{FR1 -- Multi-node autonomous operation.}
The system shall allow multiple nodes to start independently and participate in the same distributed deployment without requiring a central coordinator.

\textbf{Acceptance criterion.}
When a cluster of at least two configured nodes is started, each node becomes operational as an autonomous process and exposes its local service interfaces without depending on a single master node.

\medskip
\textbf{FR2 -- Local sensor ingestion.}
Each node shall be able to manage one or more local sensor sources or sensor connections and ingest their readings into the distributed system state.

\textbf{Acceptance criterion.}
Given at least one configured sensor on a node, the node eventually exposes fresh sensor records associated with its own origin through the read API or introspection API.

\medskip
\textbf{FR3 -- Cluster-wide sensor-state dissemination.}
The system shall propagate sensor updates among nodes through decentralized peer-to-peer replication, including push-pull anti-entropy exchanges, so that readings produced on one node become visible on other reachable nodes.

\textbf{Acceptance criterion.}
If a node generates new sensor readings while communication paths to other nodes are available, those readings eventually appear in the replicated sensor-state view exposed by the other reachable nodes.

\medskip
\textbf{FR4 -- Deterministic timestamp-based conflict resolution.}
The system shall resolve concurrent or conflicting updates using a deterministic timestamp-based merge policy (Last-Write-Wins with deterministic tie-breaking) so that replicas can converge to the same winning value for each sensor record.

\textbf{Acceptance criterion.}
When two or more updates compete for the same logical sensor record, all nodes that receive the same set of updates eventually expose the same winner for that record.

\medskip
\textbf{FR5 -- Membership discovery and peer knowledge propagation.}
The system shall allow nodes to discover peers through bootstrap and peer-to-peer membership exchange, then maintain a distributed view of cluster membership.

\textbf{Acceptance criterion.}
After startup and sufficient message exchange, a node exposes a membership view containing the peers it has discovered directly or indirectly.

\medskip
\textbf{FR6 -- Adaptive liveness observation.}
The system shall provide operators with a node-local assessment of peer liveness based on heartbeat behavior and an adaptive failure detector so that missing or degraded peers can be observed during experiments.

\textbf{Acceptance criterion.}
If a known peer becomes unreachable for long enough, at least one observing node eventually reports that peer as \texttt{suspected} or \texttt{dead} in its membership/introspection output.

\medskip
\textbf{FR7 -- Recovery after crashes or temporary disconnections.}
The system shall allow a node that restarts or reconnects after temporary unavailability to recover the replicated state needed to rejoin normal operation by synchronizing with reachable peers. Recovery is based on exchanged delta information and full-state synchronization, not on durable local persistence.

\textbf{Acceptance criterion.}
After a node crash or temporary isolation, once communication with at least one sufficiently up-to-date peer is restored, the recovering node eventually exposes a replicated state aligned with the rest of the reachable cluster.

\medskip
\textbf{FR8 -- Read-only observability interface.}
The system shall provide a read-only programmatic interface for observing sensor state, membership, topology-related information, events, and operational metrics.

\textbf{Acceptance criterion.}
A client can query documented HTTP \texttt{GET} endpoints and obtain structured snapshots of the node's current distributed-system view without modifying runtime state.

\medskip
\textbf{FR9 -- Push-pull anti-entropy exchange.}
The replication layer shall support push-pull anti-entropy exchanges between peers so that each side can both advertise local updates and request missing information.

\textbf{Acceptance criterion.}
During normal operation, reachable peers can send delta updates and request missing updates from one another, so each node can obtain sensor records it does not yet store.

\medskip
\textbf{FR10 -- Phi Accrual failure suspicion output.}
The membership subsystem shall expose liveness states derived from a Phi Accrual Failure Detector, allowing operators to distinguish peers reported as \texttt{alive}, \texttt{suspected}, or \texttt{dead}.

\textbf{Acceptance criterion.}
If heartbeat arrival patterns from a peer degrade or stop, phi-based suspicion eventually crosses the configured thresholds and the peer is reported as \texttt{suspected} or \texttt{dead} in membership/introspection output.

\medskip
\textbf{FR11 -- Full-state rejoin synchronization.}
When a node rejoins after crash, restart, or long partition, it shall support full-state synchronization from reachable peers to rebuild a coherent in-memory replica when incremental recovery is insufficient.

\textbf{Acceptance criterion.}
After prolonged unavailability, a rejoining node can request and apply a full sensor-state snapshot from a reachable peer and eventually align with the reachable cluster state, assuming such a peer still holds the relevant state.

\medskip
\textbf{FR12 -- Human-oriented monitoring support.}
The system shall support human observation of the running cluster through a dashboard or equivalent visual monitoring interface.

\textbf{Acceptance criterion.}
An observer can inspect live cluster information from a web browser through a dashboard that consumes the exposed read-only API of a node.

\subsection{Non-Functional Requirements}
\label{non-functional-requirements}

\textbf{NFR1 -- Decentralization.}
Normal operation shall not depend on any central coordinator or single control component.

\textbf{Acceptance criterion.}
The system continues to ingest local readings, maintain local state, and serve read-only observability data on surviving nodes even if one arbitrary peer becomes unavailable.

\medskip
\textbf{NFR2 -- Eventual consistency.}
The replicated sensor state shall provide eventual, not immediate, consistency across nodes.

\textbf{Acceptance criterion.}
In the absence of new conflicting updates and after communication is restored, reachable replicas converge to the same winning values for the same sensor records.

\medskip
\textbf{NFR3 -- Partition tolerance with degraded semantics.}
The system shall remain available for local operation during temporary network partitions in experimental deployments, accepting that different partitions may temporarily diverge.

\textbf{Acceptance criterion.}
During a partition, nodes in each connected component continue local ingestion and observation; after the partition heals, their replicas eventually converge again.

\medskip
\textbf{NFR4 -- Fault observability.}
The system shall make failures and recovery visible enough to support distributed-systems experimentation and debugging.

\textbf{Acceptance criterion.}
Operators can inspect status changes, replication activity, and recovery evidence through exposed membership, event, or metric views.

\medskip
\textbf{NFR5 -- Reproducibility of experiments.}
The platform shall support repeatable validation scenarios on development machines and course evaluation environments.

\textbf{Acceptance criterion.}
Using the documented setup and provided Docker Compose topologies, an operator can start a predefined cluster and reproduce representative convergence, partition, and recovery experiments without relying on external managed services.

\medskip
\textbf{NFR6 -- Extensibility of sensor sources.}
The system shall be open to additional sensor producers without changing the conceptual role of the rest of the distributed runtime. This extensibility is architectural rather than dynamic: adding a new built-in sensor type may still require code and configuration mapping in the sensor manager.

\textbf{Acceptance criterion.}
New sensor sources can be introduced as additional producers while preserving the same downstream ingestion, replication, and observation capabilities; the platform does not require a live plugin hot-swap mechanism for this requirement.

\medskip
\textbf{NFR7 -- Operational scalability at course-project scale.}
The platform shall support small and medium experimental clusters rather than only a fixed two-node demo.

\textbf{Acceptance criterion.}
The documented deployment can be instantiated with multi-node topologies, such as six-node and twelve-node clusters, and still provides replicated state and observability functions for the whole experiment.

\subsection{Implementation Constraints}
\label{implementation-constraints}

\textbf{IR1 -- Self-contained local deployment.}
The project shall be executable on commodity development machines without requiring paid cloud services, institution-managed infrastructure, or external managed services.

\textbf{Justification.}
The system is intended for course-project evaluation and distributed-systems experimentation. Its deployment model must therefore remain portable and reproducible from the repository itself.

\textbf{Acceptance criterion.}
An operator can start one or more nodes using the repository code, documented dependencies, and provided deployment artifacts without provisioning external infrastructure.

\medskip
\textbf{IR2 -- Reproducible predefined topologies.}
The project shall provide predefined multi-node deployment configurations that can be reused for demonstrations, validation, and fault-injection experiments.

\textbf{Justification.}
Validation requires repeatable cluster scenarios with stable node identities, network coordinates, sensor configurations, and observation endpoints.

\textbf{Acceptance criterion.}
An operator can launch the documented predefined topologies and obtain consistent node identities, sensor sources, peer relationships, and monitoring interfaces across repeated runs.

\medskip
\textbf{IR3 -- Course-scale experimental scope.}
The implementation shall target small and medium experimental clusters suitable for local execution and course validation, rather than production-scale IoT deployments.

\textbf{Justification.}
The system is designed to demonstrate distributed-systems properties under controlled conditions, not to satisfy industrial-scale availability, performance, or device-management requirements.

\textbf{Acceptance criterion.}
The provided configurations and validation scenarios support representative multi-node deployments suitable for evaluating convergence, partition behavior, recovery, and observability at course-project scale.

\subsection{Relevant Distributed System Features}
\label{ds-features}

Not all classical distributed-system features have the same priority in this project. Since \textit{Distributed Sensor Hub} is a decentralized monitoring and experimentation platform, some features are central design drivers, while others are intentionally secondary.

\textbf{Transparency.}
Only partially relevant. The system hides low-level transport details from observers, but it does \emph{not} aim to hide the fact that it is distributed. On the contrary, replication lag, topology evolution, liveness suspicion, and recovery dynamics are part of what the project wants to make observable. Failure transparency is therefore limited by design.

\medskip
\textbf{Fault tolerance and dependability.}
Highly relevant. Nodes may crash, restart, or become temporarily disconnected, and the platform is expected to continue operating on surviving nodes while later recovering convergence. Availability and recoverability are more important than uninterrupted global consistency. Integrity is also relevant in the sense that replicas must converge deterministically rather than drift indefinitely.

\medskip
\textbf{Scalability.}
Relevant, but at moderate scale. The project is not meant for Internet-scale deployments; however, it must scale beyond a trivial two-node example and remain meaningful on multi-node laboratory topologies. The main concern is preserving understandable behavior as the number of peers and sensor updates grows.

\medskip
\textbf{Security and trust.}
Only marginally relevant in the current scope. The platform does not manage personal data or adversarial multi-tenant access, and the main users are trusted operators and observers in controlled environments. Security hardening is therefore not a primary project driver, although it remains a possible future extension.

\medskip
\textbf{Resource sharing.}
Relevant. Nodes share distributed knowledge rather than centralized hardware resources: replicated sensor winners, membership evidence, topology views, and operational metadata. Coordination is required to keep those shared views convergent despite asynchrony and failures.

\medskip
\textbf{Openness, interoperability, and heterogeneity.}
Relevant. The project benefits from clean interfaces between sensor producers, node runtime, and observation clients. This matters because simulated sensors are only one current source of data, while the same ingestion boundary can later support heterogeneous adapters and external tools.

\medskip
\textbf{Evolvability and maintainability.}
Highly relevant. As a course project intended to demonstrate multiple distributed-systems concerns, the platform must remain understandable, modular, and extensible, especially for adding new sensors, observability surfaces, and validation scenarios.

\medskip
\textbf{Performance, concurrency, and communication efficiency.}
Relevant, but not under hard real-time constraints. The system executes many concurrent activities such as sensor production, replication, heartbeats, gossip exchange, and HTTP observation. Efficiency matters because excessive communication can obscure convergence behavior, but the project does not target strict latency guarantees.

\medskip
\textbf{Economy and costs.}
Relevant mainly as a practical constraint. The project should run on ordinary development hardware and avoid external infrastructure dependencies, which supports reproducibility, portability, and low evaluation cost.

\clearpage
